Wiske wrote a 2024-digit positive integer on the blackboard. In each round of the game she erases the last digit of the integer, let this digit be $d$, and writes down the sum of the remaining number and $2d$ in place of the old number. She repeats the same steps with the newly obtained number. After a certain number of rounds, Wiske found that the new number obtained was the same as the number in the last round and she stopped the game. What is the smallest possible 2024-digit integer that Wiske started with in this game?

(Proposed by Kai Chen)
\begin{solusi}
    \begin{remark*}[Motivasi]
        Udah, ini soal emang ajaib dan idenya didapat dari observasi banyak \textit{pattern}. Dan \textit{step} krusial lainnya adalah melihat keadaan akhir atau \textit{goal} dari bilangan ini adalah apa.
    \end{remark*}
    Misalkan bilangan di keadaan saat ini adalah $a=10x+d$ dengan $x$ adalah bilangan asli dan $d$ adalah bilangan bulat nonnegatif satu digit. Setelah melakukan satu \textit{step}, bilangan itu akan menjadi $b=x+2d$.

    \begin{lemma*}
        Bilangan pada \textit{step} terakhir adalah 19.
        \begin{proof}
            Bilangan terakhir tercapai saat $a=b$ atau
            \begin{align*}
                10x+d &= x+2d\\
                9x &= d
            \end{align*}
            Karena $d$ adalah bilangan bulat satu digit dan $x$ adalah bilangan asli, maka haruslah $x=1$ dan $d=9$. Berarti bilangan terakhir ini adalah $19$.
            
            Perhatikan jika kita sudah mendapatkan bilangan 19, maka bilangan pada \textit{step} selanjutnya adalah $1+2\times 9 = 19$, bilangan yang sama. Terbukti.
        \end{proof}
    \end{lemma*}

    Sekarang observasi modulo 19 pada bilangan tersebut.
    \begin{align*}
        2(10x+d) \equiv 19x + (x+2d) \equiv x+2d \mod 19.
    \end{align*}
    Perhatikan karena bilangan terakhir adalah 19 yang habis dibagi 19. Maka dari hubungan modulo di atas, semua bilangan sebelumnya juga habis dibagi 19 yang mudah dibuktikan dengan induksi.

    Selanjutnya mudah dilihat bahwa bilangan-bilangan itu akan membentuk barisan monoton turun karena untuk $x>1$ berlaku
    $$9x > 9 \ge \implies 10x-x > 2d - d \implies 10x+d > x+2d$$

    Oleh karena itu bilangan 2024 digit terkecil yang memenuhi adalah kelipatan 19 terkecil yang mempunyai 2024 digit. Perhatikan bahwa

    \begin{align*}
        10^4 &\equiv (10^2)^2 \equiv 100^2 \equiv 5^2 \equiv 25 \equiv 6 \mod 19\\
        \implies 10^3 &\equiv 1000 \equiv 12 \mod 19\\
        \implies 10^7 &\equiv 10^4\cdot 10^3 \equiv 6 \cdot 12 \equiv 72 \equiv 15 \mod 19.
    \end{align*}
    Lalu, berdasarkan Fermat's Little Theorem
    \begin{align*}
        10^{18} \equiv 10^{19-1} \equiv 1 \mod 19
    \end{align*}
    Sehingga
    \begin{align*}
        10^{2023} = (10^{18})^{112}10^7 \equiv 1^{112} \cdot 15 \equiv -4 \mod 19.
    \end{align*}
    Berarti bilangan 2024 digit terkecil yang habis terbagi 19 adalah $10^{2023} + 4$ yang merupakan jawaban yang kita cari.
\end{solusi}